\scanpage{313}
\mspage{\deleted{\unsure{2}}229}
% archive sheet number: 312

\diagram{\begin{tikzcd}[ampersand replacement=\&,row sep=small,column sep=small]
\&k\arrow[r,"w"]\arrow[d,no head]\&\overset{2}{\overline{k}^{\unsure{*|2}}}\arrow[d,no head]\\
i'\arrow[r,shift left]\arrow[r,shift right,no head]\&f(k)\arrow[r]\&f(\unsure{k'|\overline{k}})
\end{tikzcd}}

Mais, \unsure{en appliquant|on applique} $\mathrm{F}2$ sur $f$ à la double flèche en question \diagram{\begin{tikzcd}[ampersand replacement=\&,cramped,column sep=small] i'\arrow[r,shift left=.6ex,no head]\arrow[r,shift right=.6ex,no head]\&f(k) \end{tikzcd}}, on trouve $w:k\longrightarrow\overline{k}$ telle que

\[
\underset{\substack{\Vert\\ f(wu_1)u'}}{f(w)(f(u_1)u')}=
\underset{f(ww_1)v'}{\underbrace{f(w)(f(\unsure{w_2|w_1})v')}}\;,
\]

et remplaçant $k$ par $\overline{k}^{\unsure{*}}$, et $u_1,v_1$ par $wu_1,wv_1$, on trouve ce qu'on veut.

2) PF(2) pour $ {}_{i'}\backslash I$, i.e. on se donne un diagramme ci-contre, avec $f(u_1)u'=f(u_2)u'$, \deleted{\unsure{\(\xi\)}} \inserted{soit $v'$,} (double-flèche dans $ {}_{i'}\backslash I$ de $(i,u')$ dans $(j,v')$), \deleted{\unsure{On}} veut l'égaliser à droite. Or

\diagram{\begin{tikzcd}[ampersand replacement=\&,column sep=tiny,row sep=small]
\&i\arrow[r,shift left,"u_1"]\arrow[r,shift right,"u_2"']\arrow[d,no head]\&j\arrow[r,dashed,"w"]\arrow[d,no head]\&j'\\
i'\arrow[r,"u'"]\&f(i)\arrow[r,shift left,"f(u_1)"]\arrow[r,shift right,"f(u_2)"']\&f(j)\arrow[r,dashed,no head,"f(w')"]\&f(j')
\end{tikzcd}}

$I$ étant \unsure{filtrant|filtrante}, il existe $w:j\unsure{\mathbin{-}|\longrightarrow}j'$ telle que $wu_1=wu_2$, or cette flèche définit \deleted{\unsure{s}} une flèche $(j,v')\longrightarrow(j',f(w')v')$, qui égalise la double flèche donnée.

$b\Rightarrow b'$ est clair

$\unsure{b')|b)}\Rightarrow a)$. Comme les $ {}_{i'}\backslash I$ sont filtrantes, donc $0$-connexes, $f$ est cofinal, il reste à prouver que $I$ est filtrante, \deleted{\unsure{donc}} à vérifier \deleted{\unsure{qui est}} $\mathrm{PF}0$ ($I$ \deleted{\unsure{0}}connexe) $\mathrm{PF}1$, $\mathrm{PF}2$ pour $I$. $I$ est $0$-connexe, car $I'$ l'est \deleted{(étant filtrante)}, et $f$ cofinal i.e. $W_0$-coasphérique, d'où $\mathrm{PF}0$.

